\chapter{PS–PL հաղորդակցություն և runtime կոնֆիգուրացիա}
\label{ch:comm}

Zynq-ի մշակող համակարգը (PS) ծրագրավորվող տրամաբանության (PL) հետ կապվում է
\textbf{AXI}-ի միջոցով՝ երկու ուղիներով.
\begin{itemize}
  \item \textbf{Data Path} (AXI-Stream / FIFO) — տվյալների փոխանցում PL$\rightarrow$PS։
  \item \textbf{Control Path} (AXI4-Lite) — FPGA մոդուլների կոնֆիգուրացիա՝ ռեգիստրների միջոցով։
\end{itemize}

Web հարթակը պետք է կարողանա փոխել FPGA-ի վարքը \textbf{առանց bitstream
վերագեներացնելու}՝ AXI4-Lite ռեգիստրների միջոցով (Աղյուսակ~\ref{tab:regmap})։

\begin{table}[htbp]
  \centering
  \caption{AXI4-Lite կոնֆիգուրացիոն ռեգիստրների քարտեզի օրինակ.}
  \label{tab:regmap}
  \begin{tabular}{ll@{\hspace{2cm}}ll}
    \toprule
    \textbf{Հասցե} & \textbf{Ռեգիստր} & \textbf{Հասցե} & \textbf{Ռեգիստր} \\
    \midrule
    0x00 & CONTROL       & 0x18 & NOISE\_LEVEL   \\
    0x04 & SENSOR\_TYPE  & 0x1C & THRESHOLD      \\
    0x08 & SAMPLE\_RATE  & 0x20 & FILTER\_ENABLE \\
    0x0C & SIGNAL\_TYPE  & 0x24 & FILTER\_WINDOW \\
    0x10 & AMPLITUDE     & 0x28 & ANOMALY\_ENABLE\\
    0x14 & FREQUENCY     & 0x2C & ANOMALY\_LEVEL \\
    \bottomrule
  \end{tabular}
\end{table}

Կոնֆիգուրացիայի ուղին է՝ Web Platform $\rightarrow$ Backend/API $\rightarrow$
Ethernet $\rightarrow$ Zynq PS $\rightarrow$ AXI4-Lite $\rightarrow$ FPGA ռեգիստրներ
$\rightarrow$ RTL մոդուլներ։

\notebox{Պետք է հստակ տարբերել \textbf{runtime կոնֆիգուրացիան} (արդեն իրականացված
ապարատի պարամետրերի փոփոխություն՝ պարտադիր) և \textbf{FPGA reconfiguration}-ը
(նախապես գեներացված bitstream-ների միջև ընտրություն՝ ընդլայնում, ոչ MVP)։}

Բուֆերացումն իրականացվում է FIFO/BRAM-ով՝ Processing IP $\rightarrow$ AXI FIFO
$\rightarrow$ Zynq PS։ Բարձր տվյալահոսքի դեպքում AXI4-Stream $\rightarrow$ AXI DMA
$\rightarrow$ DDR ուղին դիտվում է որպես ընդլայնում։
